\subsection{Editing}
\label{sec:editing}

\begin{figure}[htbp]
\centering
% Width uses \textwidth; height auto-scales with !
\resizebox{\textwidth}{!}{%
\begin{tikzpicture}
% --- 1. Define leaf nodes (papers, citations added) ---
\node[leaf] (safety_papers) {
SafeEdit/DINM \citep{wang2024detoxifying}, DEPN \citep{wu2023depn}, Editing Attack \citep{chen2024canediting}, BadEdit \citep{li2024badedit}, MeGen \citep{qiu2024megen}, Concept Backdoors \citep{grimes2024concept}, EvilEdit \citep{wang2024eviledit}, Guo et al. \citep{guo2024backdoor}, Youssef et al. \citep{youssef2025position}, Uppaal et al. \citep{uppaal2024model}.
};
\node[leaf, below=1cm of safety_papers] (factuality_papers) {
IKE \citep{zheng2023can}, ROME \citep{meng2022locating}, MEMIT \citep{meng2022memit}, GRACE \citep{hartvigsen2024aging}, WISE \citep{wang2024wise}, Shi et al. \citep{shi2024retrieval}, Fei et al. \citep{fei2024retrieval}, Geva et al. \citep{geva2020transformer}, Ma et al. \citep{ma2024perturbation}, HalluEditBench \citep{huang2025halluedit}, Gu et al. \citep{gu-etal-2024-model}, Yang et al. \citep{yang-etal-2024-butterfly}.
};
\node[leaf, below=1cm of factuality_papers] (fairness_papers) {
LSDM \citep{cai2024locating}, MG-Editing \citep{qin2024mitigating}, Yan et al. \citep{yan2024potential}.
};
% --- 2. Define main sub-nodes (Trustworthiness dimensions) ---
% Positioned to the left of corresponding leaf nodes
\node[sub, left=0.5cm of safety_papers] (safety) {Safety \& Privacy};
\node[sub, left=0.5cm of factuality_papers] (factuality) {Factuality \& Faithfulness};
\node[sub, left=0.5cm of fairness_papers] (fairness) {Fairness};
% --- 3. Define root node ---
% Use the calc library to vertically center the root node relative to the sub-nodes
\node[main, left=2.5cm of {$(safety.west)!0.5!(fairness.west)$}] (editing) {Editing (§\ref{sec:editing})};
% --- 4. Draw connection lines ---
% a. Draw a main trunk from the root node to a branching point
\coordinate (fork) at ($(editing.east)+(1.25cm,0)$);
\draw[line] (editing.east) -- (fork);
% b. Connect from main branching point to each sub-node
\draw[line] (fork) |- (safety.west);
\draw[line] (fork) |- (factuality.west);
\draw[line] (fork) |- (fairness.west);
% c. Connect main sub-nodes directly to leaf nodes
\draw[line] (safety.east) -- (safety_papers.west);
\draw[line] (factuality.east) -- (factuality_papers.west);
\draw[line] (fairness.east) -- (fairness_papers.west);
\end{tikzpicture}%
}
\caption{A taxonomy of Editing methods for enhancing LLM trustworthiness, based on the dimensions, methods, and representative works discussed in Section~\ref{sec:editing}.}
\label{fig:editing_taxonomy}
\end{figure}

\subsubsection{Core Principle: Precise and Localized Parameter Updates}
% Defines Model Editing as making precise, localized modifications to an LLM's parameters to alter knowledge, aiming to avoid the catastrophic forgetting associated with full fine-tuning.



% \subsection{A Typology of Editing Mechanisms}


Knowledge editing~\citep{wang2024knowledge} is designed as an efficient
paradigm to correct erroneous or outdated factual knowledge in LLMs, considering
the high cost of retraining LLMs from scratch. Thus, it can be naturally adopted to enhance the factuality of LLMs. Recently, research has shown that knowledge editing techniques, especially the \textbf{Locate-then-edit} ones, can also be applied in other trustworthiness dimensions such as safety, privacy, and fairness, because there could be a high correlation between specific neurons in LLMs and such trustworthiness properties. 

% \subsubsection{Memory-Based Methods}
% Augmenting the model with an explicit, editable fact store.

% \subsubsection{Parameter-Based Methods}
% Directly locating and modifying the specific weights responsible for a piece of knowledge.

% \subsubsection{Editing as a Form of Unlearning}
% Framing editing techniques as a mechanism for knowledge erasure.

\subsubsection{Applications in Trustworthiness Dimensions}


\paragraph{Factuality}
% Used to correct factual inaccuracies, update outdated information, and add new knowledge directly at the source. 
Knowledge editing is mostly motivated by reducing hallucinations and enhancing factuality in LLMs. In General, knowledge editing can be categorized into four paradigms: \textbf{Fine-tuning}, \textbf{In-context editing}, \textbf{Locate-then-edit}, and \textbf{Memory-based}. The latter three ones can be regarded as inference-time methods.
% First, \textit{Fine-tuning} offers a straightforward approach of updating a model’s knowledge. Although it can be computationally costly and risks overfitting or catastrophic forgetting, techniques such as parameter-efficient tuning and early stopping help mitigate these drawbacks~\citep{gangadhar2024model,zhu2020modifying,wang2024roselora}. Second, 
First, \textit{In-context editing} is a training-free method that enables models to incorporate new knowledge directly within the input context~\citep{zheng2023can,shi2024retrieval,fei2024retrieval}. For example, \citet{zheng2023can} propose \textbf{In-Context Knowledge Editing (IKE)}, which uses demonstration-based prompting strategies to inject, generalize, and preserve facts, which shows competitive  performance compared to parameter-updating methods, achieving high efficacy while without modifying the parameters.  Second, motivated by findings that Transformer MLP layers store factual knowledge~\citep{geva2020transformer,ma2024perturbation}, \textit{Locate-then-edit} methods such as~\citep{meng2022locating,meng2022memit} aim to first identify the relevant neurons or layers and then replace them with new ones. For example, \citet{meng2022locating} investigate how GPT-like LLMs store and recall factual knowledge. They use causal tracing to show that factual recall depends heavily on MLP modules, especially at the final token of a subject in a prompt. They then introduce \textbf{Rank-One Model Editing (ROME)}, a simple algebraic method to directly insert or modify factual associations by updating weights in these MLPs. Third, \textit{Memory-based} approaches typically maintain a dedicated memory module for storing and updating knowledge~\citep{hartvigsen2024aging,wang2024wise}. For example, \textbf{GRACE}~\citep{hartvigsen2024aging} adopts a discrete codebook as the external memory while leaving the original model parameters unchanged. When faced with queries involving edited knowledge, it employs an adaptor that adjusts inter-layer transformations using values retrieved from the codebook.



\paragraph{Safety and Privacy}
% Used to neutralize harmful knowledge or detoxify behavior by editing associations leading to toxic outputs.

Knowledge Editing traditionally focuses on correcting inaccurate or outdated information within LLMs while minimizing unintended side effects on their broader outputs.  However, it also brings emerging safety risks for LLMs~\citep{youssef2025position}. Typically, \citet{chen2024canediting} reframe Knowledge Editing as a potential safety threat, introducing the concept of an \textbf{Editing Attack}, which is a type of potential malicious use of editing methods to manipulate LLMs to generate harmful content. Editing Attacks pose two primary risks: \textit{Misinformation Injection} and \textit{Bias Injection}. Misinformation Injection bypasses an LLM’s intrinsic safety alignment to implant misleading claims (e.g., “\texttt{Vaccines contain microchips}”) into models such as Llama3-8B, enabling the propagation of falsehoods through open-source model communities. Bias Injection, on the other hand, either introduces isolated biased statements (e.g., “\texttt{Poor performer in math is female}”) or undermines the overall fairness of LLMs with minimal modification (e.g., inserting a single stereotypical statement about women can cause LLMs to generate more biased answers to not only gender-related questions, but also race-related ones).

Besides harmful information injection, another form of misuse of knowledge editing involves implanting backdoors into LLMs~\citep{li2024badedit,qiu2024megen,grimes2024concept} or into text-to-image generation models~\citep{wang2024eviledit,guo2024backdoor}. The inherent properties of model editing make it feasible to embed such backdoors with relatively low cost, raising further security and reliability concerns. For example, \citet{li2024badedit} highlight the shortcomings of conventional backdoor attack approaches, which typically demand large amounts of poisoned data and significant computational overhead, often at the expense of a model’s overall functionality. To overcome these limitations, the authors introduces \textbf{BadEdit}, a lightweight editing strategy that alters only a small subset of parameters, enabling backdoor insertion with as few as 15 samples while maintaining strong performance on benign tasks. Unlike prior methods, BadEdit avoids task-specific tuning and extensive poisoned datasets; instead, it embeds trigger–target associations directly into the parameter space through key–value representation techniques. This design yields robustness to fine-tuning, near-perfect attack success rates, and minimal interference with unrelated tasks, which indicates that BadEdit could be a more practical and efficient alternative to traditional data- and weight-poisoning methods.

Apart from the potential safety challenges, knowledge editing can also be utilized to reduce model toxicity or privacy leakage~\citep{wu2023depn,wang2024detoxifying,uppaal2024model}. For example, \citet{wang2024detoxifying}  introduce \textbf{SafeEdit}, a benchmark for evaluating detoxification of LLMs via knowledge editing, covering nine unsafe categories, 48 attack prompts, and multiple evaluation metrics and \textbf{DINM} (Detoxifying with Intraoperative Neural Monitoring), a simple but effective method that locates toxic regions in LLMs through semantic differences between safe and unsafe responses and edits those parameters directly using just one adversarial instance. \citet{wu2023depn} propose a post-processing framework \textbf{DEPN} to reduce privacy leakage in large language models by identifying and editing ``privacy neurons'', which refer to neurons that store sensitive information like names or phone numbers. DEPN consists of a privacy neuron detector (using gradient integration to score neurons' contributions to private outputs), a privacy editor (zeroing top-ranked neurons), and a privacy aggregator (batch processing across sentences).


\paragraph{Fairness}
Prior works have identified that many fairness problems in LLMs are rooted from biased training data. 
As a result, many biases and stereotypes in human society are learned by models during training and therefore manifested in LLM outputs in generative tasks.
To remove such biases, one promising direction is to ``correct'' the model's social viewpoint, understanding of social norms, or language styles.
This can be achieved through model editing, which identifies and changes internal knowledge of models through various post-training and inference-time techniques.

\textbf{Locate-then-Edit for Fairness.} First proposed by \citet{meng2022locating}'s work, one of the most effective frameworks for model editing is the ``locate-then-edit'' framework.
Inspired by this method, \citet{cai2024locating} uses causal mediation analysis to identify that gender bias primarily originates from the bottom MLP modules in models. 
Based on this finding, the authors further propose the Least Square Debias Method (LSDM), a knowledge-editing technique that directly modifies the weights of these specific MLP modules to remove gender bias associated with occupational terms without retraining the model or affecting its other capabilities.
Similarly, \citet{qin2024mitigating} proposes MG-Editing to mitigate gender bias in code LLMs.
MG-Editing operates at five levels of granularity—from the full model down to individual neurons—to first locate and then directly adjust the specific parameters responsible for bias, ensuring that models' code generation capabilities remain intact while removing biases.

\textbf{General Insights}
More generally, \citet{yan2024potential} benchmarks 7 existing model editing algorithms for the task of social debiasing in LLMs. 
The authors found that while current methods can effectively mitigate specific biases without harming the model's general knowledge, they struggle with generalizing the debiasing effect to semantically similar sentences and performing sequential edits.

% \begin{itemize}
%     % \item \citep{cai2024locating}
%     % \item \citep{yan2024potential}
%     \item \citep{qin2024mitigating}
% \end{itemize}
% % An emerging application focused on directly editing stereotypical associations stored in the model's weights.

% \subsubsection{Privacy}
% Used to remove PII by targeting and modifying the parameters or "privacy neurons" that store personal data.



\subsubsection{Discussion: The Challenge of Generalization, Locality and Scalability}

The primary strength of knowledge editing methods lies in their  computational efficiency and plug-in nature. These approaches are designed to make targeted updates to a model's knowledge without the need for retraining the whole parameters, which can otherwise be computationally expensive and time-consuming. Because knowledge editing methods mostly operate in a post-hoc and plug-in fashion, they can be easily integrated into existing architectures, applied selectively to specific tasks, and reused across different contexts.

However, recent works also shed light on the challenges of knowledge editing techniques on \textit{generalization}, \textit{locality} and \textit{scalability}~\citep{huang2025halluedit,hsueh2024editing,gu-etal-2024-model,hoelscher2023detecting,hua2024propagation,yang-etal-2024-butterfly,yang2024fall,li2023unveiling,cohen2024evaluating}. For example, with a newly developed benchmark, \textbf{HalluEditBench}, which spans 9 domains, 26 topics, and over 6,000 verified hallucinations, \citet{huang2025halluedit} show that the actual ability of knowledge editing techniques to correct real-world hallucinations may fall significantly short of what their performance on previous datasets indicates. Also, most editing methods have an unsatisfying performance on generalization and locality scores. For generalization, even when an edit fixes the original hallucinated answer, it often fails to transfer to rephrased, yes/no, multi-choice, or reversed forms of the same question, and in some cases (notably GRACE and LoRA) edits even have lower generalization scores compared to the unedited model. For locality, most methods disrupt unrelated knowledge, with overall scores under 40\% in many domains. Also, \citet{gu-etal-2024-model} provide a critical insight that even a single model edit, which may alter less than 1\% of parameters, can drastically degrade the general abilities of LLMs, with LLaMA-1 (7B) in some cases collapsing to near-zero performance across tasks. This reveals that LLMs are highly fragile to weight perturbations, contradicting the assumption that localized edits can be made harmlessly. As for sequential editing, which applies multiple edits in succession, each new edit accumulates distortions in the model’s weight space, further eroding abilities like reasoning, summarization, and question answering. 
% \citet{yang-etal-2024-butterfly} discovered that even one single edit could cause a significant performance degradation in mainstream benchmarks. 


% Strengths: Computationally efficient and precise post-hoc correction.
% Limitations: The core challenge is ensuring locality (an edit does not affect unrelated knowledge), scaling to millions of facts, and the difficulty of editing complex behaviors versus simple facts.





